% main_report74.tex
% SAMBP Technical Report TR-74/2028
% Negative-Sequence Overcurrent for IBR Networks
\documentclass[11pt, a4paper]{article}

\usepackage[T1]{fontenc}
\usepackage[utf8]{inputenc}
\usepackage[margin=25mm]{geometry}
\usepackage{amsmath, amssymb, amsthm}
\usepackage{booktabs, array, multirow, longtable}
\usepackage{graphicx}
\usepackage[table]{xcolor}
\usepackage[hidelinks]{hyperref}
\usepackage{pgfplots}
\pgfplotsset{compat=1.18}
\usepackage{tikz}
\usetikzlibrary{arrows.meta, positioning, calc, fit, backgrounds,
                shapes.geometric, shapes.misc}
\usepackage{caption}
\usepackage{subcaption}
\usepackage{parskip}
\usepackage{siunitx}
\usepackage{pifont}
\usepackage{cite}

\definecolor{sambpblue}{RGB}{0,70,140}
\definecolor{okgreen}{RGB}{0,130,60}
\definecolor{alertred}{RGB}{180,0,0}
\definecolor{warnoran}{RGB}{200,100,0}

\newcommand{\pu}{\,\text{pu}}
\newcommand{\ms}{\,\text{ms}}
\newcommand{\cmark}{\ding{51}}
\newcommand{\xmark}{\ding{55}}
\newcommand{\kibr}{k_{\mathrm{ibr}}}
\newcommand{\fgfm}{f_{\mathrm{gfm}}}

\newtheorem{finding}{Finding}
\newtheorem{remark}{Remark}

\title{\textbf{\textsf{\color{sambpblue}SAMBP Technical Report TR-74/2028}}\\[6pt]
  \Large Negative-Sequence Overcurrent Protection for IBR Networks\\[4pt]
  \normalsize Relay~46 Adaptation, GFM Misoperation Risk,
  Adaptive Polarisation: 90-Scenario Study}
\author{SAMBP Research Group\\[2pt]
  \small Smart Adaptive Multi-Bus Protection Research, IIT Madras}
\date{April 2028\quad\textbar\quad Chapter~3 (Protection Challenges)\quad\textbar\quad
  Prerequisite: TR-69 (Type-4 Wind), TR-56 (DFIG)}

\begin{document}
\maketitle
\tableofcontents
\vspace{6pt}\hrule\vspace{10pt}

\section{Background and Objectives}
\label{sec:background}

\subsection{Motivation}

Negative-sequence overcurrent (Relay~46) is a standard protection function
for detecting unbalanced faults (SLG, LL, DLG) in conventional power systems.
The fundamental assumption is that \emph{negative-sequence current flows from
source to fault}: sources have a low negative-sequence impedance and the fault
is the negative-sequence sink. This assumption breaks down in IBR networks:

\begin{itemize}
  \item \textbf{GFL IBR:} Injects near-zero negative-sequence current (converter
    maintains balanced three-phase output). Relay~46 sees reduced $I_2$ and
    may fail to operate (under-detection).
  \item \textbf{GFM IBR:} Can inject \emph{reverse} negative-sequence current
    to actively balance the network voltage (VSM virtual rotor dynamics).
    This can cause Relay~46 at adjacent buses to see $I_2$ in the wrong
    direction and misoperate.
\end{itemize}

\subsection{Objectives}

\begin{enumerate}
  \item Model negative-sequence current injection for GFL and GFM converters.
  \item Quantify Relay~46 misoperation risk for three IBR mix levels.
  \item Develop an adaptive polarisation scheme using positive-sequence voltage.
  \item Validate 90 scenarios across $3$ IBR mix levels $\times$ $30$ fault types.
\end{enumerate}

\section{Theory and Methodology}
\label{sec:theory}

\subsection{GFL Negative-Sequence Model}

A GFL converter with balanced current control injects:
$\mathbf{I}_{\mathrm{GFL}} = [I_1, 0, 0]^T$ (sequence components),
i.e., $I_2^{\mathrm{GFL}} \approx 0$.
The actual $I_2$ depends on control loop bandwidth and grid impedance asymmetry:
\begin{equation}
  I_2^{\mathrm{GFL}} = \frac{V_2}{Z_2^{\mathrm{ctrl}} + Z_{\mathrm{grid,2}}}
  \label{eq:gfl_i2}
\end{equation}
where $Z_2^{\mathrm{ctrl}} = R_{\mathrm{ctrl}} + j\omega_0 L_{\mathrm{ctrl}}$
is the effective negative-sequence control impedance ($\approx 10$--$50\pu$
for fast current controllers). The effective $I_2^{\mathrm{GFL}}$ is
$\leq 0.03\pu$ for most fault conditions.

\subsection{GFM Negative-Sequence Model}

A GFM converter with VSM control actively suppresses negative-sequence voltage:
\begin{equation}
  \mathbf{V}_{\mathrm{GFM,ref}} = V_1 e^{j\theta} + V_2^* e^{-j\theta}
    \cdot G_{\mathrm{NS}}(s)
  \label{eq:gfm_ns}
\end{equation}
where $G_{\mathrm{NS}}(s)$ is a negative-sequence suppression filter.
The resulting $I_2^{\mathrm{GFM}}$ flows \emph{out} of the GFM source
(toward the fault), which can reverse the apparent direction of $I_2$
seen at adjacent Relay~46 elements.

The reverse negative-sequence index is:
$\xi_2 = \angle(I_2) - \angle(V_2) + 180^\circ$;
$\xi_2 \in (-90^\circ, +90^\circ)$ for forward, outside for reverse.

\subsection{Adaptive Polarisation Scheme}

The proposed adaptive scheme uses the positive-sequence voltage $V_1$ as
the polarisation reference for Relay~46:
\begin{equation}
  I_{2,\mathrm{pol}} = I_2 \cdot e^{j(\angle V_1 + \phi_{\mathrm{comp}})}
  \label{eq:pol}
\end{equation}
where $\phi_{\mathrm{comp}} = \angle(Z_1/Z_2)$ compensates for the
impedance angle difference between the positive and negative sequence networks.
This polarisation is robust to GFM reverse injection because $V_1$ is
maintained during faults by the GFM source.

Pick-up criterion: $\mathrm{Re}[I_{2,\mathrm{pol}}] > I_{2,\mathrm{set}}$
AND $|I_2| > I_{2,\mathrm{min}} = 0.05\pu$.

\section{Simulation Results}
\label{sec:results}

\subsection{90-Scenario Results}

\begin{table}[ht]
\centering
\caption{TR-74 90-scenario results by IBR mix level.}
\label{tab:results}
\begin{tabular}{lccccc}
\toprule
IBR mix ($\fgfm$) & $\kibr$ & Scenarios & AGREE & Agree Rate & Status \\
\midrule
Low ($\fgfm=0.1$)    & 0.3 & 30 & 30 & 100\%  & \cmark \\
Medium ($\fgfm=0.5$) & 0.6 & 30 & 28 & 93.3\% & \cmark \\
High ($\fgfm=1.0$)   & 1.0 & 30 & 27 & 90.0\% & \cmark \\
\midrule
\textbf{Total} & --- & \textbf{90} & \textbf{85} & \textbf{94.4\%}
  & \textbf{\color{okgreen}PASS ($\geq 94\%$)} \\
\bottomrule
\end{tabular}
\end{table}

\subsection{Misoperation Risk Analysis}

\begin{table}[ht]
\centering
\caption{Relay~46 misoperation rate: standard vs.\ adaptive polarisation.}
\label{tab:misop}
\begin{tabular}{lccc}
\toprule
IBR mix & Standard 46 misop.\ rate & Adaptive 46 misop.\ rate & Reduction \\
\midrule
Low ($\fgfm=0.1$)    & 0\%   & 0\%   & --- \\
Medium ($\fgfm=0.5$) & 13\%  & 6.7\% & 48\% \\
High ($\fgfm=1.0$)   & 27\%  & 10\%  & 63\% \\
\bottomrule
\end{tabular}
\end{table}

\begin{finding}[GFM injection reverses $I_2$ direction]
At $\fgfm = 1.0$, 27\% of standard Relay~46 elements misoperate due to GFM
reverse negative-sequence injection. The adaptive positive-sequence polarisation
scheme reduces misoperation to 10\%, a 63\% improvement.
\end{finding}

\begin{finding}[Residual misoperation at high GFM mix]
Even with adaptive polarisation, 10\% misoperation rate at $\fgfm = 1.0$
indicates that Relay~46 is not fully reliable in all-GFM networks.
High-impedance fault (HIF) scenarios are the primary failure mode:
the $|I_2| < I_{2,\mathrm{min}} = 0.05\pu$ threshold is not exceeded
for HIF at $\fgfm = 1.0$. A lower threshold or auxiliary HIF detector
(as in TR-50) is recommended.
\end{finding}

\section{Conclusion}
\label{sec:conclusion}

TR-74 demonstrates that Relay~46 requires significant adaptation for IBR-rich
networks. The adaptive positive-sequence polarisation scheme reduces misoperation
by 63\% at high GFM penetration. The 90-scenario campaign achieves 94.4\%
agreement across all IBR mix levels, meeting the $\geq 94\%$ criterion.
Residual misoperation at all-GFM networks motivates further development in TR-77.

\begin{thebibliography}{99}
\bibitem{Hooshyar2019}
A.~Hooshyar and R.~Iravani,
``Microgrid protection,''
\emph{Proc.\ IEEE}, vol.~105, no.~7, pp.~1332--1353, 2017.

\bibitem{Plet2012}
C.~A.~Plet \emph{et al.},
``Fault models of inverter-interfaced distributed generators:
  experimental verification and application to fault analysis,''
in \emph{Proc.\ IEEE PES General Meeting}, San Diego, CA, 2012.

\bibitem{Nabe1986}
A.~Nabe, I.~Takahashi, and H.~Akagi,
``A new neutral-point-clamped PWM inverter,''
\emph{IEEE Trans.\ Ind.\ Appl.}, vol.~IA-17, no.~5, pp.~518--523, 1981.

\bibitem{Fortescue1918}
C.~L.~Fortescue,
``Method of symmetrical coordinates applied to the solution of polyphase
networks,''
\emph{Trans.\ AIEE}, vol.~37, pp.~1027--1140, 1918.

\bibitem{IEEE_C37_110}
IEEE Std C37.110-2007,
\emph{IEEE Guide for the Application of Current Transformers Used for
Protective Relaying Purposes}, IEEE, 2007.
\end{thebibliography}

\end{document}
